% =====================================================================
% polish_preamble.tex -- Team F talk preamble: palette, hairlines, theorem style, sizes
% (body 10pt, frame titles 12pt, tables 9pt, captions 8pt), old-style figures in text,
% lining figures in math, quiet section navigation in the footer.
% =====================================================================
\usefonttheme{serif}
\usefonttheme{professionalfonts}
\usepackage{amsmath}
\usepackage{fontspec}
\usepackage[math-style=TeX,bold-style=TeX]{unicode-math}
\setmainfont{TeX Gyre Pagella}[Numbers={Lining,Proportional}]   % Palatino text, lining figures
\setsansfont{TeX Gyre Pagella}[Numbers={Lining,Proportional}]   % no stray sans anywhere
\setmonofont{TeX Gyre Pagella}[Numbers=Lining]
\setmathfont{texgyrepagella-math.otf}                              % matching Palatino math
\usepackage{array}
\linespread{1.18}

% ---------- palette --------------------------------------------------------
\definecolor{ink}{HTML}{1F2933}
\definecolor{accent}{HTML}{23426B}
\definecolor{muted}{HTML}{4B5563}   % secondary text, lighter than ink
\definecolor{hairline}{HTML}{C9CDD2}
\definecolor{note}{HTML}{374151}     % page-bottom notes and captions (darker than muted)
\definecolor{navcur}{HTML}{C2642B}   % current section in the footer nav (muted burnt orange)
\definecolor{hlY}{HTML}{FFEB99}
\definecolor{hlG}{HTML}{CDEBC5}
\definecolor{hlR}{HTML}{F9D3C4}
\definecolor{boxred}{HTML}{B03A2E}
\definecolor{lineblue}{HTML}{2F5D8C}
\setbeamercolor{normal text}{fg=ink,bg=white}
\setbeamercolor{background canvas}{bg=white}
\setbeamercolor{structure}{fg=accent}
\setbeamercolor{alerted text}{fg=boxred}
\setbeamercolor{frametitle}{fg=accent,bg=}
\setbeamercolor{footline}{fg=accent}   % footer in the structure navy
\setbeamercolor{itemize item}{fg=muted}
\setbeamercolor{enumerate item}{fg=muted}

% ---------- sizes -------------------------------------------------------------
\setbeamerfont{frametitle}{size=\fontsize{12}{14.5}\selectfont}
\setbeamerfont{footline}{size=\fontsize{7.5}{9}\selectfont}

% ---------- layout --------------------------------------------------------------
\setbeamersize{text margin left=8mm,text margin right=8mm}
\setbeamertemplate{navigation symbols}{}
\setbeamertemplate{headline}{}
\setbeamertemplate{frametitle}{%
  \vskip 1.6mm% fixed-height title box: rule at the same y on every page
  \hbox to\textwidth{\usebeamerfont{frametitle}\usebeamercolor[fg]{frametitle}\rule{0pt}{4.2mm}\smash{\insertframetitle}\hfil}%
  \nointerlineskip\vskip 3.45mm\hbox{\color{hairline}\rule{\textwidth}{0.4pt}}\vskip -2.6mm}
% quiet agenda in the footer: the current part in accent, the others in muted small caps
\newcommand{\navitem}[2]{\ifnum\value{section}=#1{\color{navcur}\textsc{#2}}\else\textsc{#2}\fi}   % current part in burnt orange
\newcommand{\navsep}{\hspace{0.9em}}
\setbeamertemplate{footline}{%
  \usebeamerfont{footline}\usebeamercolor[fg]{footline}%
  \hbox to \paperwidth{\hskip 8mm\rlap{\textsc{Mixed boundary data}\ \textperiodcentered\ \textsc{Team F}}\hfill
    \navitem{1}{model}\navsep\navitem{2}{analysis}\navsep\navitem{3}{experiments}\navsep
    \navitem{4}{close}\hfill
    \llap{\insertframenumber\,/\,\inserttotalframenumber}\hskip 8mm}\vskip 3.6mm}
\addtobeamertemplate{frame begin}{}{\raggedright\hyphenpenalty=10000\exhyphenpenalty=10000}
\hyphenpenalty=10000 \exhyphenpenalty=10000
\setbeamertemplate{itemize item}{\raisebox{0.12ex}{\textendash}}
\setbeamertemplate{enumerate item}{\insertenumlabel.}
\setlength{\leftmargini}{1.4em}

% ---------- theorem-like: small-caps label, no coloured blocks ------------------
\setbeamertemplate{theorems}[normal font]
\setbeamertemplate{theorem begin}{\par\addvspace{\blockskip}\noindent
  {\upshape\scshape\color{accent}\inserttheoremname}%
  \ifx\inserttheoremaddition\empty\else\ {\color{muted}(\inserttheoremaddition)}\fi.\ \itshape}
\setbeamertemplate{theorem end}{\par\addvspace{\blockskip}}
\newtheorem*{proposition*}{Proposition}
\newtheorem*{remark*}{Remark}

% ---------- no emphasis marks: no boxes, underlines or highlights -------------------
\newcommand{\dfn}[1]{{\color{accent}#1}}
\newcommand{\notesize}{\fontsize{8.5}{10.6}\selectfont}         % notes/captions 8.5pt
\newcommand{\notesizesmall}{\fontsize{7.5}{9.4}\selectfont}    % AI-use block 7.5pt
\newcommand{\aside}[1]{{\color{note}\notesize\upshape #1}}
\newcommand{\figcap}[1]{{\color{note}\notesize\upshape #1\par}}   % \par inside: footnotesize leading

% ---------- tables and display spacing ---------------------------------------------
\usepackage{booktabs}
\setlength{\aboverulesep}{0.3ex}\setlength{\belowrulesep}{0.45ex}
\setlength{\heavyrulewidth}{0.6pt}\setlength{\lightrulewidth}{0.35pt}\setlength{\cmidrulewidth}{0.35pt}
\renewcommand{\arraystretch}{1.1}
\newcommand{\tabfont}{\small\addfontfeatures{Numbers={Lining,Monospaced}}}   % 9pt, lining figures so table columns align
\newcommand{\stack}[1]{\begin{tabular}[t]{@{}l@{}}#1\end{tabular}}   % manual line break inside a cell
\makeatletter
\g@addto@macro\normalsize{\setlength\abovedisplayskip{5pt}\setlength\belowdisplayskip{5pt}%
  \setlength\abovedisplayshortskip{4pt}\setlength\belowdisplayshortskip{4pt}}
\makeatother

% ---------- layout rhythm ---------------------------------------------------
% every frame body starts \bodytop below the title rule; footnote-style asides sit \asidegap above the footer
\newlength{\bodyskip}\setlength{\bodyskip}{0.6mm}
\newlength{\blockskip}\setlength{\blockskip}{3.0mm}% the one fixed gap between body blocks
\newcommand{\bodytop}{\vspace*{\bodyskip}\vspace*{1.6mm}\vspace{0pt plus 0.3fill}}
\newcommand{\tabletop}{\vspace*{\bodyskip}\vspace*{2.5mm}}   % table top rule level with the cap height of text pages
\newcommand{\closing}[1]{\par\addvspace{\blockskip}{\color{note}\notesize\upshape #1\par}\vspace{0pt plus 1fill}}
\newcommand{\figend}[1]{\par\vspace{\blockskip}{\color{note}\notesize\upshape #1\par}\vspace{0pt plus 1fill}}   % figure frames: tighter bottom gap
\newcommand{\blockgap}{\par\addvspace{\blockskip}}     % between body blocks (table -> text, list -> text)
\newcommand{\colgap}{\par\vspace{\blockskip}}       % between paragraphs inside the text column of figure frames
\newcommand{\lead}[1]{{\upshape\scshape\color{accent}#1}}   % small-caps run-in label (Goal, Remark)

% ---------- stroke weight: text drawn fill+stroke with a hairline stroke in the text colour ----
% Palatino stays Palatino; glyphs gain ~0.1pt per side, so the deck reads better when projected.
% Text remains real text (selectable, searchable); colours unchanged (xcolor sets fill and stroke).
\usepackage{pdfrender}
\newcommand{\strokeweight}{0.15pt}
\AtBeginDocument{\pdfrender{TextRenderingMode=FillStroke,LineWidth=\strokeweight,LineJoinStyle=1}}
% ---------- Project 2 additions ----------
\definecolor{todo}{HTML}{B03A2E}
\newcommand{\TODO}[1]{{\color{todo}[TODO: #1]}}

\setlength{\parskip}{0pt}
\AtBeginDocument{\renewcommand{\setminus}{\mathbin{\backslash}}}
